\section{System Overview}
\label{sec:system-overview}

This section describes \vannak{}'s architecture, component inventory, the two
event planes in detail, the end-to-end data flow, and the separation between
the hot path and the agreement path.

\subsection{Architecture}

\vannak{} ingests two related event streams, keeps recent operational state in
memory, writes raw history and metadata outbox entries to append-only segments,
delivers durable metadata to \ipto{}, and reserves \raft{} for cluster
coordination.

\begin{lstlisting}[numbers=none]
Durga pipeline runtimes
    -> typed process events + data-individual metadata events
    -> Vannak ingest API
    -> Sitas shard-local hot indexes
    -> metadata correlation and Ipto placement
    -> durable metadata outbox
    -> Ipto PostgreSQL instance selected by DataIndividualShardId
    -> query API / snapshots / alerting
    -> append-only event segments
    -> Raft-backed cluster control state
\end{lstlisting}

\subsection{Component Inventory}

\begin{description}
\item[Ingest] validates \texttt{PipelineEvent} values and rejects empty
  required fields.
\item[Hot index] \texttt{HotIndex} reduces process events into current
  per-instance state and builds metadata-impact indexes.
\item[Placement] \texttt{IptoPlacementMap} resolves a
  \texttt{DataIndividualShardId} to an \texttt{IptoInstanceId}.
\item[Mapping] \texttt{IptoMapping} translates metadata field names to \ipto{}
  attribute names.
\item[Outbox] \texttt{MetadataOutbox} and \texttt{SegmentBackedMetadataOutbox}
  hold pending durable writes.
\item[Writer] the \texttt{IptoWriter} trait and \texttt{IptoRepoWriter}
  deliver payloads idempotently into \ipto{}.
\item[Storage] \texttt{SegmentWriter} / \texttt{SegmentReader} provide
  append-only durability.
\item[Cluster] \texttt{ClusterControlState} holds compact \raft{}-replicable
  control state.
\end{description}

\subsection{Two Event Planes}

\subsubsection{Process Events}

The \texttt{PipelineEvent} type is the validated process event accepted by the
hot path. Each event carries a coarse \texttt{EventStatus} and a precise
\texttt{EventKind}.

\texttt{EventStatus} variants: \texttt{Started}, \texttt{Completed},
\texttt{Failed}, \texttt{Escalated}, \texttt{Cancelled}.

\texttt{EventKind} variants: \texttt{ProcessStarted}, \texttt{ActivityEntered},
\texttt{ActivityCompleted}, \texttt{ActivityEscalated},
\texttt{ActivityCancelled}, \texttt{GatewayTaken}, \texttt{ProcessCompleted},
\texttt{ProcessFailed}. Each kind has an \texttt{inferred\_status()} mapping to
an \texttt{EventStatus}.

\begin{lstlisting}
pub struct PipelineEvent {
    event_id: EventId,
    source_id: SourceId,
    source_sequence: SourceSequence,
    tenant_id: TenantId,
    environment_id: EnvironmentId,
    pipeline_id: PipelineId,
    process_definition_id: ProcessDefinitionId,
    process_instance_id: ProcessInstanceId,
    process_version: Option<ProcessVersion>,
    activity_id: Option<ActivityId>,
    token_id: Option<TokenId>,
    correlation_id: Option<CorrelationId>,
    business_key: Option<BusinessKey>,
    timestamp: EventTimestamp,
    status: EventStatus,
    kind: EventKind,
    error: Option<ErrorInfo>,
    metadata_refs: Vec<MetadataRef>,
    metadata_version: Option<MetadataVersion>,
    causal_parent: Option<EventId>,
    payload: Option<String>,
}
\end{lstlisting}

\subsubsection{Data-Individual Metadata Events}

The \texttt{DataIndividualMetadataEvent} type is a provenance fact for one
flowing data item. It distinguishes \texttt{PassiveMetadata} (observed at
receive/create boundaries) from \texttt{ActiveMetadata} (produced by
transformations, masking, validation, enrichment). Both are maps of
\texttt{MetadataFieldName} to \texttt{MetadataValue}.

\texttt{MetadataValue} variants: \texttt{String}, \texttt{Integer},
\texttt{Boolean}, \texttt{Timestamp}, \texttt{StringList}.

The \texttt{MetadataOperation} enum classifies the provenance operation:

\begin{lstlisting}
pub enum MetadataOperation {
    Created,
    Received,
    Transformed { plugin_name: Option<PluginName>,
                  plugin_version: Option<PluginVersion> },
    Masked { fields: Vec<String> },
    Validated { passed: bool },
    Enriched { source: Option<String> },
    Routed,
    Persisted,
}
\end{lstlisting}

Each event carries a \texttt{DataIndividualShardId} (the domain placement key)
and an \texttt{IdempotencyKey} derived from the data-individual id and the
metadata-event id.

\subsection{Data Flow}

The full path of a captured metadata event is:

\begin{enumerate}
\item \textbf{Ingest} a process event or metadata event at the typed API.
\item \textbf{Validate} required fields via \texttt{PipelineEvent::validate}.
\item \textbf{Route} the event to the owning \sitas{} shard.
\item \textbf{Hot index} reduces process state and indexes metadata impact.
\item \textbf{Metadata event} carries passive and active metadata.
\item \textbf{Resolve shard} maps \texttt{DataIndividualShardId} to an
  \texttt{IptoInstanceId} through the placement map.
\item \textbf{Map attributes} translates fields to \ipto{} attributes via
  \texttt{IptoWritePayload::from\_event}.
\item \textbf{Outbox} persists the encoded payload to a segment before it is
  visible as pending.
\item \textbf{IptoWriter} delivers the payload idempotently.
\item \textbf{\ipto{}} stores the durable provenance unit.
\end{enumerate}

\subsection{Hot Path versus Agreement Path}

The hot path is shard-local and append-only:

\begin{lstlisting}[numbers=none]
event -> validate -> route -> shard-local reducer/index -> segment writer
\end{lstlisting}

The agreement path uses \raft{} only for decisions that must be agreed across
the cluster:

\begin{lstlisting}[numbers=none]
ownership decision  -> Raft log -> committed cluster state
checkpoint decision -> Raft log -> committed manifest
sealed segment      -> Raft log -> committed segment manifest index
\end{lstlisting}

\begin{vnote}
High-volume event ingest must never be routed through \raft{}. Consensus is
reserved for membership, ownership, placement maps, metadata-version
activation, checkpoint manifests, sealed-segment manifests, and writer leases.
\end{vnote}

\subsection{Feature Flag Matrix}

\begin{table}[htbp]
\centering
\begin{tabular}{lll}
\hline
Capability & Required feature & Key types \\
\hline
Core domain & (none) & \texttt{HotIndex}, \texttt{MetadataOutbox} \\
\ipto{} writes & \texttt{ipto-writer} & \texttt{IptoRepoWriter} \\
State machine & \texttt{raft} & \texttt{ClusterStateMachine} \\
Node binary & \texttt{node} & \texttt{vannak-node} \\
\hline
\end{tabular}
\caption{Capabilities enabled by each feature flag.}
\end{table}

\subsection{Crate Layout}

The crate root re-exports the public types of every module. The dependency-free
modules (\texttt{ingest}, \texttt{process}, \texttt{data}, \texttt{metadata},
\texttt{durga}, \texttt{index}, \texttt{query}, \texttt{ipto},
\texttt{storage}, \texttt{cluster}, \texttt{observability}, \texttt{runtime})
compile with no external crates. The \texttt{ipto\_adapter} and
\texttt{raft\_sm} modules are compiled only when their feature flags are
enabled.

\endinput
